\section{Problem Restatement}

The problem asks for more than an evaluation ranking. At the surface level, the candidate suppliers must be compared under several indicators such as cost, reliability, and effective loss. However, the final research question is operational: after the candidate set is identified, an executable supply plan must still be generated under demand and feasibility constraints. Therefore, the paper must answer both the screening question and the planning question in one continuous route.

For paper organization, the task is decomposed into three connected subquestions. The first subquestion is how to construct a credible indicator system and normalize heterogeneous supplier information. The second subquestion is how to use the selected candidate set to produce a period-by-period allocation or order plan. The third subquestion is how to judge whether the recommendation remains acceptable when the external environment, especially loss or demand pressure, changes. This decomposition is important because it determines the later section rhythm: the model is not a single block, but a chain from evaluation to decision to validation.

The core deliverables of the paper are therefore not limited to one score table. A complete draft should contain at least four visible outputs: a supplier-comparison artifact, a candidate-selection result, an executable plan table, and a feasibility or scenario-check artifact. Only when these outputs appear together can the paper plausibly support a final recommendation rather than a descriptive ranking.

In research use, the value of this route lies in traceability. The evaluation layer must explain why certain suppliers remain plausible, the planning layer must explain how those suppliers are used in an executable decision, and the validation layer must explain why the resulting plan is credible enough for later review or discussion. This is the problem restatement that governs the rest of the paper.
